\documentclass[11pt]{article}
\usepackage[margin=1in]{geometry}
\usepackage{booktabs}
\usepackage{caption}
\usepackage{amsmath}
\usepackage{float}
\usepackage{graphicx}
\usepackage{adjustbox}

\begin{document}

\section*{Asymmetric Demand Multipliers by Shock Size: v3}

This note reports the v3 pooled positive/negative specification. Unlike v2,
which estimates positive- and negative-shock subsamples separately, v3 keeps
both signs in the same monthly cross-section and estimates six sign-specific
bin slopes in one regression.

For each demand type and month, the pooled shock-size bins are inherited from
the pooled regression input file. Let $d_{n,t}$ denote demand and let
$d^{(k)}_{n,t}$ be the pooled-$\sigma$ bin variables:
\[
d^{(1)}_{n,t} = d_{n,t}\mathbf{1}\{|d_{n,t}| < \sigma_t\}, \quad
d^{(2)}_{n,t} = d_{n,t}\mathbf{1}\{|d_{n,t}| \in [\sigma_t,2\sigma_t]\}, \quad
d^{(3)}_{n,t} = d_{n,t}\mathbf{1}\{|d_{n,t}| > 2\sigma_t\}.
\]
The six regressors are then
\[
d^{(k),+}_{n,t} = d^{(k)}_{n,t}\mathbf{1}\{d_{n,t} > 0\}, \qquad
d^{(k),-}_{n,t} = d^{(k)}_{n,t}\mathbf{1}\{d_{n,t} < 0\},
\]
for $k \in \{1,2,3\}$. The monthly Fama--MacBeth regression is
\[
r_{n,t} = \alpha_t
        + \sum_{k=1}^3 M_k^+ d^{(k),+}_{n,t}
        + \sum_{k=1}^3 M_k^- d^{(k),-}_{n,t}
        + \Gamma X_{n,t} + \varepsilon_{n,t}.
\]

The regression is estimated separately for BMI, FIT, and OFI demand. Columns
(1)--(3) use BMI, columns (4)--(6) use FIT, and columns (7)--(9) use OFI.
Within each demand type, the three columns correspond to no controls,
predictor controls, and predictor plus liquidity controls. BMI specifications
also include the BMI-specific control set used in the main price-impact tables.

\paragraph{Implementation detail.}
If a monthly cross-section omits a coefficient because the corresponding bin
has no observations, v3 keeps that monthly coefficient missing when computing
Fama--MacBeth averages. It does not replace omitted coefficients with zero.
This differs from the Stata \texttt{statsby \_b[var]} display behavior, which
can effectively report omitted coefficients as zero.

The difference rows are Fama--MacBeth averages of monthly paired differences,
not differences of the printed coefficient averages. These can diverge when
one bin has omitted monthly coefficients. For example, in FIT negative column
(3), bin 2 averages 2.17 over the months in which bin 2 is estimated, while
bin 3 averages 2.25 over the months in which bin 3 is estimated. The
``Bin 3 - Bin 2'' row averages the month-by-month difference only over months
where both coefficients are estimated; in that common-month sample the mean
bin 2 coefficient is 1.98 and the mean bin 3 coefficient is 2.25, giving a
paired difference of 0.27.

\paragraph{Takeaway.}
The positive-demand domain shows strong concavity: multipliers are largest in
the smallest shock bin and smallest in the largest shock bin for FIT and OFI.
This is not consistent with a short-sales-constraint interpretation under
which the multiplier should rise with shock size. On the negative-demand side,
OFI also shows strong concavity. FIT is weaker on the negative side: the
largest-bin multiplier is below the smallest-bin multiplier, but the adjacent
large-minus-medium difference is small and imprecise. BMI remains noisy because
the BMI sample has only 21 monthly cross-sections. Overall, v3 does not show
evidence that $M$ increases with shock size, so the evidence is not consistent
with leverage constraints generating larger multipliers for larger shocks.

\begin{table}[H]
\centering
\caption{v3 shock-size bin summary statistics}
\label{tab:v3_bin_summary}
\begin{tabular}{llrrrr}
\toprule
Shock & Type & Bin & Obs & Mean $d_{n,t}$ & SD $d_{n,t}$ \\
\midrule
Negative & BMI & 1 & 3,922 & -0.001889 & 0.001698 \\
Negative & BMI & 2 & 472 & -0.011188 & 0.003155 \\
Negative & BMI & 3 & 287 & -0.020162 & 0.006973 \\
Negative & FIT & 1 & 176,791 & -0.001351 & 0.001312 \\
Negative & FIT & 2 & 40,711 & -0.005666 & 0.002131 \\
Negative & FIT & 3 & 17,203 & -0.011023 & 0.004662 \\
Negative & OFI & 1 & 247,374 & -0.013672 & 0.011270 \\
Negative & OFI & 2 & 40,829 & -0.056677 & 0.017312 \\
Negative & OFI & 3 & 18,668 & -0.137339 & 0.070365 \\
\addlinespace
Positive & BMI & 1 & 3,973 & 0.002133 & 0.001837 \\
Positive & BMI & 2 & 654 & 0.010492 & 0.003105 \\
Positive & BMI & 3 & 392 & 0.020672 & 0.006904 \\
Positive & FIT & 1 & 233,213 & 0.001444 & 0.001496 \\
Positive & FIT & 2 & 48,811 & 0.007076 & 0.002600 \\
Positive & FIT & 3 & 25,849 & 0.015667 & 0.008279 \\
Positive & OFI & 1 & 187,161 & 0.012776 & 0.011188 \\
Positive & OFI & 2 & 24,051 & 0.056741 & 0.017554 \\
Positive & OFI & 3 & 9,661 & 0.134889 & 0.054641 \\
\bottomrule
\end{tabular}
\end{table}

\begin{table}[H]
\centering
\caption{v3 pooled six-variable asymmetric multipliers by shock-size bin}
\label{tab:v3_asym_multipliers}
\begin{adjustbox}{max width=\textwidth}
\begin{tabular}{lccccccccc}

  \hline \multicolumn{10}{c}{Panel A: Positive shocks} \\

  \hline

  & \multicolumn{3}{c}{BMI} & \multicolumn{3}{c}{FIT} & \multicolumn{3}{c}{OFI} \\

  \cmidrule(lr){2-4} \cmidrule(lr){5-7} \cmidrule(lr){8-10}

  & (1) & (2) & (3) & (1) & (2) & (3) & (1) & (2) & (3) \\

  Bin 1 & $1.54^{*}$ & $1.77^{*}$ & $2.05^{**}$ & $9.52^{***}$ & $7.53^{***}$ & $6.95^{***}$ & $3.19^{***}$ & $3.26^{***}$ & $3.66^{***}$ \\

   & $(0.89)$ & $(0.92)$ & $(0.83)$ & $(1.38)$ & $(1.09)$ & $(0.86)$ & $(0.22)$ & $(0.18)$ & $(0.16)$ \\

  Bin 2 & $0.53$ & $0.39$ & $0.63$ & $6.42^{***}$ & $5.51^{***}$ & $5.39^{***}$ & $2.22^{***}$ & $2.26^{***}$ & $2.49^{***}$ \\

   & $(0.65)$ & $(0.70)$ & $(0.67)$ & $(0.88)$ & $(0.72)$ & $(0.64)$ & $(0.18)$ & $(0.16)$ & $(0.14)$ \\

  Bin 3 & $0.82^{**}$ & $0.57$ & $0.62$ & $4.04^{***}$ & $3.64^{***}$ & $3.65^{***}$ & $1.57^{***}$ & $1.57^{***}$ & $1.69^{***}$ \\

   & $(0.40)$ & $(0.37)$ & $(0.39)$ & $(0.58)$ & $(0.53)$ & $(0.52)$ & $(0.17)$ & $(0.16)$ & $(0.15)$ \\

  Bin 2 - Bin 1 & $-1.01$ & $-1.38$ & $-1.42$ & $-3.10^{***}$ & $-2.02^{***}$ & $-1.56^{***}$ & $-0.98^{***}$ & $-1.00^{***}$ & $-1.17^{***}$ \\

   & $(1.17)$ & $(1.08)$ & $(0.99)$ & $(0.76)$ & $(0.64)$ & $(0.56)$ & $(0.11)$ & $(0.10)$ & $(0.11)$ \\

  Bin 3 - Bin 2 & $0.29$ & $0.19$ & $-0.00$ & $-2.38^{***}$ & $-1.87^{***}$ & $-1.74^{***}$ & $-0.65^{***}$ & $-0.70^{***}$ & $-0.80^{***}$ \\

   & $(0.59)$ & $(0.69)$ & $(0.65)$ & $(0.70)$ & $(0.58)$ & $(0.57)$ & $(0.12)$ & $(0.11)$ & $(0.11)$ \\

  Bin 3 - Bin 1 & $-0.72$ & $-1.19$ & $-1.42^{*}$ & $-5.49^{***}$ & $-3.89^{***}$ & $-3.31^{***}$ & $-1.62^{***}$ & $-1.70^{***}$ & $-1.97^{***}$ \\

   & $(0.91)$ & $(0.88)$ & $(0.81)$ & $(1.24)$ & $(1.00)$ & $(0.84)$ & $(0.16)$ & $(0.14)$ & $(0.14)$ \\

  \hline \multicolumn{10}{c}{Panel B: Negative shocks} \\

  \hline

  & \multicolumn{3}{c}{BMI} & \multicolumn{3}{c}{FIT} & \multicolumn{3}{c}{OFI} \\

  \cmidrule(lr){2-4} \cmidrule(lr){5-7} \cmidrule(lr){8-10}

  & (1) & (2) & (3) & (1) & (2) & (3) & (1) & (2) & (3) \\

  Bin 1 & $4.13^{**}$ & $2.19^{*}$ & $1.81^{*}$ & $4.49^{**}$ & $4.04^{***}$ & $3.58^{***}$ & $2.51^{***}$ & $2.52^{***}$ & $2.67^{***}$ \\

   & $(1.64)$ & $(1.16)$ & $(1.02)$ & $(1.84)$ & $(1.38)$ & $(1.08)$ & $(0.15)$ & $(0.11)$ & $(0.11)$ \\

  Bin 2 & $2.58^{***}$ & $1.77^{***}$ & $1.60^{***}$ & $2.42^{**}$ & $2.21^{***}$ & $2.17^{***}$ & $1.65^{***}$ & $1.65^{***}$ & $1.71^{***}$ \\

   & $(0.84)$ & $(0.61)$ & $(0.57)$ & $(1.01)$ & $(0.73)$ & $(0.62)$ & $(0.12)$ & $(0.09)$ & $(0.09)$ \\

  Bin 3 & $1.03^{**}$ & $0.79^{**}$ & $0.77^{**}$ & $2.33^{***}$ & $2.27^{***}$ & $2.25^{***}$ & $0.56^{***}$ & $0.57^{***}$ & $0.60^{***}$ \\

   & $(0.41)$ & $(0.40)$ & $(0.38)$ & $(0.69)$ & $(0.58)$ & $(0.56)$ & $(0.07)$ & $(0.06)$ & $(0.06)$ \\

  Bin 2 - Bin 1 & $-1.55$ & $-0.43$ & $-0.21$ & $-1.84^{*}$ & $-1.67^{*}$ & $-1.32^{*}$ & $-0.86^{***}$ & $-0.87^{***}$ & $-0.96^{***}$ \\

   & $(1.09)$ & $(0.86)$ & $(0.76)$ & $(1.11)$ & $(0.94)$ & $(0.79)$ & $(0.07)$ & $(0.07)$ & $(0.07)$ \\

  Bin 3 - Bin 2 & $-1.55$ & $-0.98$ & $-0.83$ & $0.14$ & $0.28$ & $0.27$ & $-1.10^{***}$ & $-1.08^{***}$ & $-1.11^{***}$ \\

   & $(1.03)$ & $(0.79)$ & $(0.73)$ & $(0.87)$ & $(0.75)$ & $(0.71)$ & $(0.09)$ & $(0.08)$ & $(0.07)$ \\

  Bin 3 - Bin 1 & $-3.10^{*}$ & $-1.40$ & $-1.04$ & $-1.43$ & $-1.16$ & $-0.88$ & $-1.95^{***}$ & $-1.95^{***}$ & $-2.07^{***}$ \\

   & $(1.83)$ & $(1.32)$ & $(1.17)$ & $(1.60)$ & $(1.22)$ & $(1.00)$ & $(0.11)$ & $(0.10)$ & $(0.10)$ \\

  \hline

  Predictor controls & No & Yes & Yes & No & Yes & Yes & No & Yes & Yes \\

  Liquidity controls & No & No & Yes & No & No & Yes & No & No & Yes \\

  Observations & 9,908 & 9,908 & 9,908 & 542,578 & 542,578 & 542,578 & 527,744 & 527,744 & 527,744 \\

  $R^2$ & 0.061 & 0.150 & 0.177 & 0.013 & 0.067 & 0.082 & 0.063 & 0.113 & 0.134 \\

  Marginal $R^2$ & 0.029 & 0.022 & 0.021 & 0.013 & 0.008 & 0.007 & 0.063 & 0.054 & 0.055 \\

  \hline

\end{tabular}


\end{adjustbox}
\end{table}

\paragraph{Balanced-month variant.}
The table below repeats the v3 pooled six-variable specification after
dropping demand-type/month cells where any of the six sign-bin regressors has
zero nonzero observations. This restriction drops no BMI months and no OFI
months. It drops four FIT months: 200106, 200203, 200403, and 201303. In each
dropped FIT month, the negative largest-shock bin is empty; in 200106 and
200203 the negative medium-shock bin is also empty. The balanced FIT sample
therefore uses 115 monthly cross-sections instead of 119.

\begin{table}[H]
\centering
\caption{v3 pooled six-variable asymmetric multipliers, balanced sign-bin months}
\label{tab:v3_asym_multipliers_balanced}
\begin{adjustbox}{max width=\textwidth}
\begin{tabular}{lccccccccc}

  \hline \multicolumn{10}{c}{Panel A: Positive shocks} \\

  \hline

  & \multicolumn{3}{c}{BMI} & \multicolumn{3}{c}{FIT} & \multicolumn{3}{c}{OFI} \\

  \cmidrule(lr){2-4} \cmidrule(lr){5-7} \cmidrule(lr){8-10}

  & (1) & (2) & (3) & (1) & (2) & (3) & (1) & (2) & (3) \\

  Bin 1 & $1.54^{*}$ & $1.77^{*}$ & $2.05^{**}$ & $9.93^{***}$ & $7.83^{***}$ & $7.15^{***}$ & $3.19^{***}$ & $3.26^{***}$ & $3.66^{***}$ \\

   & $(0.89)$ & $(0.92)$ & $(0.83)$ & $(1.40)$ & $(1.11)$ & $(0.88)$ & $(0.22)$ & $(0.18)$ & $(0.16)$ \\

  Bin 2 & $0.53$ & $0.39$ & $0.63$ & $6.64^{***}$ & $5.69^{***}$ & $5.52^{***}$ & $2.22^{***}$ & $2.26^{***}$ & $2.49^{***}$ \\

   & $(0.65)$ & $(0.70)$ & $(0.67)$ & $(0.91)$ & $(0.74)$ & $(0.66)$ & $(0.18)$ & $(0.16)$ & $(0.14)$ \\

  Bin 3 & $0.82^{**}$ & $0.57$ & $0.62$ & $4.15^{***}$ & $3.75^{***}$ & $3.74^{***}$ & $1.57^{***}$ & $1.57^{***}$ & $1.69^{***}$ \\

   & $(0.40)$ & $(0.37)$ & $(0.39)$ & $(0.59)$ & $(0.55)$ & $(0.54)$ & $(0.17)$ & $(0.16)$ & $(0.15)$ \\

  Bin 2 - Bin 1 & $-1.01$ & $-1.38$ & $-1.42$ & $-3.29^{***}$ & $-2.14^{***}$ & $-1.62^{***}$ & $-0.98^{***}$ & $-1.00^{***}$ & $-1.17^{***}$ \\

   & $(1.17)$ & $(1.08)$ & $(0.99)$ & $(0.78)$ & $(0.66)$ & $(0.58)$ & $(0.11)$ & $(0.10)$ & $(0.11)$ \\

  Bin 3 - Bin 2 & $0.29$ & $0.19$ & $-0.00$ & $-2.49^{***}$ & $-1.94^{***}$ & $-1.78^{***}$ & $-0.65^{***}$ & $-0.70^{***}$ & $-0.80^{***}$ \\

   & $(0.59)$ & $(0.69)$ & $(0.65)$ & $(0.72)$ & $(0.60)$ & $(0.59)$ & $(0.12)$ & $(0.11)$ & $(0.11)$ \\

  Bin 3 - Bin 1 & $-0.72$ & $-1.19$ & $-1.42^{*}$ & $-5.77^{***}$ & $-4.09^{***}$ & $-3.41^{***}$ & $-1.62^{***}$ & $-1.70^{***}$ & $-1.97^{***}$ \\

   & $(0.91)$ & $(0.88)$ & $(0.81)$ & $(1.27)$ & $(1.02)$ & $(0.87)$ & $(0.16)$ & $(0.14)$ & $(0.14)$ \\

  \hline \multicolumn{10}{c}{Panel B: Negative shocks} \\

  \hline

  & \multicolumn{3}{c}{BMI} & \multicolumn{3}{c}{FIT} & \multicolumn{3}{c}{OFI} \\

  \cmidrule(lr){2-4} \cmidrule(lr){5-7} \cmidrule(lr){8-10}

  & (1) & (2) & (3) & (1) & (2) & (3) & (1) & (2) & (3) \\

  Bin 1 & $4.13^{**}$ & $2.19^{*}$ & $1.81^{*}$ & $3.76^{**}$ & $3.43^{**}$ & $3.12^{***}$ & $2.51^{***}$ & $2.52^{***}$ & $2.67^{***}$ \\

   & $(1.64)$ & $(1.16)$ & $(1.02)$ & $(1.82)$ & $(1.36)$ & $(1.06)$ & $(0.15)$ & $(0.11)$ & $(0.11)$ \\

  Bin 2 & $2.58^{***}$ & $1.77^{***}$ & $1.60^{***}$ & $2.19^{**}$ & $1.99^{***}$ & $1.98^{***}$ & $1.65^{***}$ & $1.65^{***}$ & $1.71^{***}$ \\

   & $(0.84)$ & $(0.61)$ & $(0.57)$ & $(1.01)$ & $(0.73)$ & $(0.61)$ & $(0.12)$ & $(0.09)$ & $(0.09)$ \\

  Bin 3 & $1.03^{**}$ & $0.79^{**}$ & $0.77^{**}$ & $2.33^{***}$ & $2.27^{***}$ & $2.25^{***}$ & $0.56^{***}$ & $0.57^{***}$ & $0.60^{***}$ \\

   & $(0.41)$ & $(0.40)$ & $(0.38)$ & $(0.69)$ & $(0.58)$ & $(0.56)$ & $(0.07)$ & $(0.06)$ & $(0.06)$ \\

  Bin 2 - Bin 1 & $-1.55$ & $-0.43$ & $-0.21$ & $-1.57$ & $-1.44$ & $-1.15$ & $-0.86^{***}$ & $-0.87^{***}$ & $-0.96^{***}$ \\

   & $(1.09)$ & $(0.86)$ & $(0.76)$ & $(1.11)$ & $(0.93)$ & $(0.78)$ & $(0.07)$ & $(0.07)$ & $(0.07)$ \\

  Bin 3 - Bin 2 & $-1.55$ & $-0.98$ & $-0.83$ & $0.14$ & $0.28$ & $0.27$ & $-1.10^{***}$ & $-1.08^{***}$ & $-1.11^{***}$ \\

   & $(1.03)$ & $(0.79)$ & $(0.73)$ & $(0.87)$ & $(0.75)$ & $(0.71)$ & $(0.09)$ & $(0.08)$ & $(0.07)$ \\

  Bin 3 - Bin 1 & $-3.10^{*}$ & $-1.40$ & $-1.04$ & $-1.43$ & $-1.16$ & $-0.88$ & $-1.95^{***}$ & $-1.95^{***}$ & $-2.07^{***}$ \\

   & $(1.83)$ & $(1.32)$ & $(1.17)$ & $(1.60)$ & $(1.22)$ & $(1.00)$ & $(0.11)$ & $(0.10)$ & $(0.10)$ \\

  \hline

  Predictor controls & No & Yes & Yes & No & Yes & Yes & No & Yes & Yes \\

  Liquidity controls & No & No & Yes & No & No & Yes & No & No & Yes \\

  Observations & 9,908 & 9,908 & 9,908 & 523,528 & 523,528 & 523,528 & 527,744 & 527,744 & 527,744 \\

  $R^2$ & 0.061 & 0.150 & 0.177 & 0.013 & 0.068 & 0.083 & 0.063 & 0.113 & 0.134 \\

  Marginal $R^2$ & 0.029 & 0.022 & 0.021 & 0.013 & 0.008 & 0.007 & 0.063 & 0.054 & 0.055 \\

  \hline

\end{tabular}


\end{adjustbox}
\end{table}

In the most controlled balanced FIT specification, positive multipliers are
7.15, 5.52, and 3.74 across bins 1--3. Negative multipliers are 3.12, 1.98,
and 2.25. The negative-side largest-minus-smallest difference remains -0.88
and imprecise, while the medium-to-large difference is 0.27 and imprecise.
The balanced restriction therefore resolves the differing-month-set issue in
the original v3 presentation but does not change the substantive conclusion:
FIT is strongly concave on the positive side and weak/flat on the negative
side, while OFI remains strongly concave on both sides.

\paragraph{Controlled-specification highlights.}
In the most controlled specification, positive FIT multipliers fall from 6.95
in the smallest bin to 3.65 in the largest bin; positive OFI multipliers fall
from 3.66 to 1.69. Negative OFI multipliers fall from 2.67 to 0.60, with a
largest-minus-smallest difference of -2.07. Negative FIT falls from 3.58 to
2.25 between the smallest and largest bins, but the difference is imprecise.
BMI coefficients are not the main evidence because the sample has only 21
months and the estimates are noisy.

\end{document}
